
La démonstration du
théorème 3.6 occupera la fin de cet exposé. Donnons quelques préliminaires.

\textbf{Proposition 3.7.} \uline{Quels que soient} $x\in \Gr K^{\bullet}(X)_{\mathbb Q}$, $y\in \Gr K^{\bullet}(Y)_{\mathbb Q}$, \uline{on a} :
\[
 f_{\mathrm{gr}}\bigl(x\cdot f^{\mathrm{gr}}(y)\bigr)
 =
 f_{\mathrm{gr}}(x)\,y .
\]
Par additivité, on peut supposer \(x\) et \(y\) homogènes; soit $x\in \Gr^iK^{\bullet}(X)_{\mathbb Q}$, $y\in \Gr^jK^{\bullet}(Y)_{\mathbb Q}$.
Soient $x'\in \Fil^iK^{\bullet}(X)_{\mathbb Q}$, $y'\in \Fil^jK^{\bullet}(Y)_{\mathbb Q}$ des éléments relevant \(x\) et \(y\). Alors la formule de projection donne :
\[
 f_*\bigl(x'\cdot f^*(y')\bigr)=f_*(x')\,y'.
\]
D'après 3.2, cette égalité s'écrit dans
\(\Fil^{i+j-d}K^{\bullet}(Y)_{\mathbb Q}\). On en déduit une égalité dans
\(\Gr^{i+j-d}K^{\bullet}(Y)_{\mathbb Q}\), et on vérifie immédiatement que c'est bien l'égalité cherchée.

\textbf{Lemme 3.8.} \uline{Soient} \(f:X\to Y\), \uline{et} \(g:Y\to Z\) \uline{deux morphismes vérifiant les hypothèses de} 3.6. \uline{Alors} \(g\circ f\) \uline{les vérifie également}. \uline{Si} \(x\) \uline{est un élément de} \(K^{\bullet}(X)\), \uline{et si} 3.6 \uline{est vérifié par} \(f\) \uline{pour l'élément} \(x\), \uline{et par} \(g\) \uline{pour l'élément} \(f_*(x)\), \uline{alors il est vérifié par} \(g\circ f\) \uline{pour l'élément} \(x\).

Par hypothèse, on a :
\begin{equation}\tag{3.8.1}
 \ch(f_*(x))=
 f_{\mathrm{gr}}\bigl(\ch(x)\cdot\Todd(\check T_f)\bigr),
\end{equation}
et :
\begin{equation}\tag{3.8.2}
 \ch(g_*(f_*(x)))=
 g_{\mathrm{gr}}\bigl(\ch(f_*(x))\cdot\Todd(\check T_g)\bigr).
\end{equation}

En combinant (3.8.1) et (3.8.2), on obtient :
\[
\ch(g_*(f_*(x)))=
g_{\mathrm{gr}}\!\left(
 f_{\mathrm{gr}}\bigl(\ch(x)\Todd(\check T_f)\bigr)\cdot \Todd(\check T_g)
\right).
\]
